\documentclass[11pt,a4paper]{article}

\usepackage{lmodern}
\usepackage[T1]{fontenc}
\usepackage[utf8]{inputenc}
\usepackage[margin=2.6cm]{geometry}
\usepackage{amsmath}
\usepackage{amssymb}
\usepackage{booktabs}
\usepackage{microtype}
\usepackage{xcolor}
\usepackage[round]{natbib}
\usepackage[hidelinks]{hyperref}

\newcommand{\code}[1]{\texttt{#1}}

% A numbered, ruled algorithm block in base LaTeX. The `algorithm` and
% `algorithmicx` packages are not in the TeX Live subset the paper workflow
% installs (texlive-latex-extra ships neither; they sit in texlive-science),
% so the environment is defined here rather than pulled in.
\newcounter{algo}
\newenvironment{algobox}[2]{%
  \par\addvspace{\medskipamount}%
  \refstepcounter{algo}\label{#2}%
  \noindent\rule{\linewidth}{0.8pt}\par\nobreak\vspace{2pt}%
  \noindent\textbf{Algorithm \thealgo.} #1\par\nobreak\vspace{2pt}%
  \noindent\rule{\linewidth}{0.4pt}\par\nobreak\vspace{2pt}%
  \begin{list}{\arabic{algoline}:}{%
    \usecounter{algoline}%
    \setlength{\leftmargin}{2.2em}\setlength{\labelwidth}{1.6em}%
    \setlength{\itemsep}{1pt}\setlength{\parsep}{0pt}\setlength{\topsep}{2pt}}%
}{%
  \end{list}\nobreak\vspace{2pt}%
  \noindent\rule{\linewidth}{0.8pt}\par\addvspace{\medskipamount}%
}
\newcounter{algoline}
\DeclareMathOperator{\diag}{diag}
\DeclareMathOperator*{\argmin}{arg\,min}

\title{\code{pyhrp}: Hierarchical Risk Parity and Schur Complementary
Allocation on an Explicit Cluster Tree}

\author{Thomas Schmelzer\thanks{\url{https://github.com/tschm/pyhrp}}}

\date{\today}

\begin{document}

\maketitle

\begin{abstract}
\code{pyhrp} is a small Python package that implements Hierarchical Risk Parity
(HRP) and its Schur-complement extension by allocating weights on an
\emph{explicit}, inspectable binary tree rather than by returning a bare weight
vector. This paper documents the construction the package actually performs ---
the correlation-to-distance map, the linkage and bisection trees, the recursive
inverse-variance split, and the Schur augmentation of the two child blocks ---
and reports what that construction produces on the twenty-asset panel shipped
with the test suite. Of three claims the package documents about its Schur
allocator, one holds exactly and two do not. The one that holds is worth
relying on: at $\gamma=0$ the Schur allocator reproduces HRP bit for bit, so it
is a safe drop-in. The two that fail are stated without qualification in the
package's own documentation --- that $\gamma=1$ recovers the global
minimum-variance portfolio, and that variance falls monotonically in $\gamma$ ---
and we give a two-asset counterexample to the first and a $76/200$ failure rate
for the second. The gap is in the split rule, which stays
inverse-variance at every $\gamma$; we make its consequences precise, since a
user choosing $\gamma$ is choosing between two behaviours that the
documentation currently conflates.
\end{abstract}

\section{Introduction}

Mean-variance optimisation \citep{markowitz1952} inverts a covariance matrix.
That inversion is the source of both its appeal and its notorious instability:
the smallest eigenvalues of an estimated covariance matrix are the least well
determined, and inversion weights them most heavily, so small changes in the
input produce large changes in the output \citep{michaud1989}. The standard
responses are to improve the estimate --- shrinkage \citep{ledoit2004} --- or to
abandon the optimisation --- equal weights \citep{demiguel2009}.

Hierarchical Risk Parity \citep{lopezdeprado2016} takes a third route. It never
inverts anything. Instead it uses the correlation matrix twice, and only in ways
that are robust to estimation noise: once to build a hierarchy over the assets
by agglomerative clustering, and once to compute sub-portfolio variances, which
are quadratic forms rather than solves. Weights fall out of a recursive top-down
split of a unit budget between the two children of every node, inversely
proportional to each child's variance. The result is long-only by construction,
sums to one without a constraint being imposed, and requires no expected returns.

The idea that a correlation matrix over financial assets carries hierarchical
structure predates HRP \citep{mantegna1999}; HRP's contribution is to make that
structure the allocation mechanism rather than a visualisation.

\subsection{What this package adds}

\code{pyhrp} makes the tree a first-class object. \code{hrp()} returns the root
\code{Cluster} of the weighted tree, not a weight series, so the intermediate
allocation at every node remains available for inspection: which cluster
received what fraction of the budget, and how that fraction was subdivided. The
weight vector is one attribute of the root, recovered as
\code{root.portfolio.weights}.

Three allocators share that tree and the same input contract:

\begin{itemize}
  \item \code{risk\_parity} --- HRP as published: split by the plain block
        variance of each child.
  \item \code{schur\_risk\_parity} --- split by a Schur-augmented block variance
        \citep{cotton2024}, interpolating with a parameter $\gamma \in [0,1]$.
  \item \code{one\_over\_n} --- a hierarchical equal-weight allocation,
        exposed as a generator that yields one complete portfolio per tree
        level, so the allocation can be watched as the hierarchy deepens.
\end{itemize}

Section~\ref{sec:method} states the construction, Section~\ref{sec:impl}
describes the implementation choices that are not implied by the mathematics,
and Section~\ref{sec:results} reports numerical results, including the two
documentation claims that do not hold.

\section{Method}
\label{sec:method}

Let $R \in \mathbb{R}^{T \times N}$ be a matrix of simple returns for $N$ assets
over $T$ periods, with sample covariance $\Sigma \in \mathbb{R}^{N \times N}$
and sample correlation $C \in \mathbb{R}^{N \times N}$.

\subsection{From correlation to a tree}

Correlations are turned into distances by the standard map
\begin{equation}
  d_{ij} = \sqrt{\frac{1 - C_{ij}}{2}},
  \label{eq:dist}
\end{equation}
clipped to $[0,1]$ before the square root and with an exactly zero diagonal.
Perfectly correlated assets sit at distance $0$, uncorrelated assets at
$1/\sqrt{2}$, perfectly anticorrelated assets at $1$.

The condensed form of $d$ is passed to SciPy's agglomerative clustering
\citep{virtanen2020} under one of four linkage methods --- \code{single},
\code{complete}, \code{average} or \code{ward} \citep{ward1963}, the default.
The resulting linkage matrix is converted into a binary tree of \code{Cluster}
nodes, each leaf carrying the column index of one asset.

The correlation matrix is validated before use. Fewer than two assets, or any
non-finite entry, raises. A non-finite diagonal entry is reported with the
offending asset named, because the cause is almost always a constant price
series, whose zero variance makes its correlations undefined.

\paragraph{The bisection variant.}
Setting \code{bisection=True} keeps only the \emph{leaf order} that the chosen
linkage induces and discards the tree shape. The ordered list of leaves is then
split in half, recursively, until singletons remain. This is closer to the
construction described in the original HRP paper, and it produces a balanced
tree by definition: depth is $\lceil \log_2 N \rceil$ regardless of how
chain-like the linkage was.

A bisection tree has no linkage matrix of its own, so one is rebuilt for
plotting. The merge height assigned to an internal node is its \emph{cluster
diameter},
\begin{equation}
  h(v) = \max_{i,j \in \mathrm{leaves}(v)} d_{ij},
\end{equation}
which is monotone along the tree --- a parent's leaf set contains each child's,
so its maximum can only grow --- and monotonicity is exactly what SciPy requires
of that column.

\subsection{The recursive split}

Every node holds a \code{Portfolio}: a mapping from asset name to weight. A leaf
holds its own asset at weight $1$. An internal node combines its children.

For an internal node $v$ with children $\ell$ and $r$, let $w_\ell$ and $w_r$
denote the child portfolios --- each already a unit-sum allocation over its own
leaf set --- and let $\Sigma_\ell$, $\Sigma_r$ be the corresponding diagonal
blocks of $\Sigma$. Define the child variances
\begin{equation}
  v_\ell = w_\ell^{\top} \Sigma_\ell w_\ell,
  \qquad
  v_r = w_r^{\top} \Sigma_r w_r .
  \label{eq:blockvar}
\end{equation}
The budget of $v$ is split
\begin{equation}
  \alpha_\ell = \frac{v_r}{v_\ell + v_r},
  \qquad
  \alpha_r = 1 - \alpha_\ell = \frac{v_\ell}{v_\ell + v_r},
  \label{eq:split}
\end{equation}
so that $\alpha_\ell v_\ell = \alpha_r v_r$: the two children contribute equally
to the parent, up to the cross term, which this rule ignores. If both variances
vanish --- riskless sub-portfolios --- the split is $\alpha_\ell = \alpha_r =
1/2$ rather than undefined. The node's portfolio is
\begin{equation}
  w_v = \alpha_\ell \, w_\ell \oplus \alpha_r \, w_r ,
\end{equation}
where $\oplus$ concatenates over disjoint asset sets. Since $\alpha_\ell +
\alpha_r = 1$ and both are non-negative, the invariant that a node's weights are
non-negative and sum to one propagates from the leaves to the root.

\begin{algobox}{Allocation on a cluster tree}{alg:alloc}
\item[] \textbf{Input:} root $\rho$, covariance $\Sigma$, per-node variance rule $\nu$.
\item $\mathcal{O} \gets$ nodes of $\rho$ in pre-order, collected with an explicit
      stack so that every parent precedes its children.
\item For each $v \in \mathrm{reverse}(\mathcal{O})$:
      \begin{itemize}
        \item if $v$ is a leaf, set $w_v \gets \{\, \mathrm{asset}(v) : 1 \,\}$;
        \item otherwise compute $(v_\ell, v_r) \gets \nu(\ell, r, \Sigma)$ ---
              Eq.~\eqref{eq:blockvar} for HRP, Eq.~\eqref{eq:schur} for Schur ---
              and set $w_v \gets \alpha_\ell w_\ell \oplus \alpha_r w_r$ by
              Eq.~\eqref{eq:split}.
      \end{itemize}
\item Return $\rho$.
\end{algobox}

Algorithm~\ref{alg:alloc} is the whole allocator. Reversing a pre-order walk
guarantees that both children carry a portfolio before their parent needs one,
which is what makes a single bottom-up pass sufficient.

\subsection{Schur complementary allocation}

Equation~\eqref{eq:blockvar} uses only the diagonal blocks of $\Sigma$. The
cross-block covariance between the two children is discarded --- at every node,
at every level. Schur Complementary Allocation \citep{cotton2024} restores part
of it.

Write the covariance restricted to the union of the two children's assets in
block form,
\begin{equation}
  \Sigma_v =
  \begin{pmatrix}
    A & B \\
    B^{\top} & D
  \end{pmatrix},
\end{equation}
with $A$ the left block and $D$ the right. The augmented blocks are
\begin{equation}
  A_\gamma = A - \gamma \, B D^{-1} B^{\top},
  \qquad
  D_\gamma = D - \gamma \, B^{\top} A^{-1} B,
  \label{eq:aug}
\end{equation}
and the split then uses
\begin{equation}
  v_\ell = w_\ell^{\top} A_\gamma \, w_\ell,
  \qquad
  v_r = w_r^{\top} D_\gamma \, w_r .
  \label{eq:schur}
\end{equation}
At $\gamma=0$ this is Eq.~\eqref{eq:blockvar} unchanged. At $\gamma=1$ the
augmented blocks are the Schur complements of $D$ and $A$ in $\Sigma_v$, each
group's covariance \emph{conditioned} on the other. Intermediate $\gamma$
interpolates linearly in the correction term.

Both augmentations require a solve against a covariance block. Collinear assets
make those blocks singular, so the implementation falls back from
\code{np.linalg.solve} to the minimum-norm least-squares solution when the solve
raises, which keeps Eq.~\eqref{eq:aug} defined rather than propagating an
exception out of a recursion.

Note what is \emph{not} changed: the split rule. Equation~\eqref{eq:split}
remains inverse-variance for every $\gamma$. Section~\ref{sec:claims} shows why
that matters.

\subsection{Hierarchical \texorpdfstring{$1/N$}{1/N}}

The third allocator ignores $\Sigma$. At level $n$ of the tree it distributes a
budget $2^{-n}$ equally among the leaves of each node at that level, accumulates
into a running portfolio, and yields a copy. A leaf that terminates early keeps
the weight it was given, so each yielded portfolio is a complete allocation over
all assets and the sequence shows how the equal-weight allocation refines as the
hierarchy is descended.

\section{Implementation}
\label{sec:impl}

The package is about 1{,}200 lines of Python over eight modules, dependent only
on NumPy \citep{harris2020}, SciPy \citep{virtanen2020}, Polars \citep{polars}
and Plotly. Covariance and correlation frames are Polars DataFrames whose
columns name the assets; the allocators translate names to indices once and
then work on NumPy arrays.

\paragraph{Iterative traversal, not recursion.}
Both the SciPy-tree conversion and the allocation walk are written as explicit
two-pass stack loops. This is not stylistic. Single linkage chains: on noisy
data its tree depth grows with $N$ rather than with $\log N$, and the recursive
forms of these two functions raised \code{RecursionError} on universes that the
mathematics handles without complaint. The iterative form removes the asset
count as a limit.

\paragraph{The allocator contract.}
\code{risk\_parity} and \code{schur\_risk\_parity} write \emph{into} the tree
they are given and return the same root. Every node's portfolio is replaced
rather than accumulated into, which makes them idempotent: re-running on an
already weighted tree --- with a different covariance matrix, or a different
$\gamma$ --- gives the same answer as running on a fresh one. The cost is that
the caller's tree is modified. \code{Dendrogram} is a frozen dataclass, but the
\code{Cluster} it holds is not, so passing \code{dendrogram.root} to an
allocator rewrites the portfolios inside that dendrogram; a caller who needs the
unweighted tree must deep-copy it. \code{one\_over\_n} differs on both counts:
it accumulates into a local buffer, leaving the input untouched, and yields
rather than returns.

\paragraph{Numerical conventions.}
Returns are simple returns from prices, with leading all-null rows dropped and
remaining nulls and NaNs filled with zero. Covariance and correlation are the
plain sample estimators (\code{np.cov}, \code{np.corrcoef}). No shrinkage is
applied anywhere: the package's robustness claim rests on avoiding the
inversion, not on regularising the estimate.

\section{Numerical results}
\label{sec:results}

All numbers below come from the twenty-asset daily price panel shipped in the
test suite as \code{tests/resources/stock\_prices.csv} ($320$ observations,
$20$ assets). Variances are of daily returns; annualised volatility multiplies
by $\sqrt{252}$. We report the effective number of positions
$N_{\mathrm{eff}} = 1 / \sum_i w_i^2$ as a concentration measure --- it equals
$N$ for equal weights and $1$ for a single position.

\subsection{Tree construction}

\begin{table}[t]
\centering
\caption{HRP allocations on the twenty-asset panel, by linkage method and tree
construction. \emph{Bisect} rebuilds the tree by halving the linkage leaf order.}
\label{tab:linkage}
\begin{tabular}{llrrrr}
\toprule
Linkage & Bisect & Variance & Ann.\ vol. & $N_{\mathrm{eff}}$ & $\max_i w_i$ \\
\midrule
\code{single}   & no  & $6.388 \times 10^{-5}$ & $0.1269$ & $7.82$  & $0.2245$ \\
\code{single}   & yes & $5.325 \times 10^{-5}$ & $0.1158$ & $13.63$ & $0.1308$ \\
\code{complete} & no  & $5.255 \times 10^{-5}$ & $0.1151$ & $11.82$ & $0.1591$ \\
\code{complete} & yes & $5.131 \times 10^{-5}$ & $0.1137$ & $13.00$ & $0.1288$ \\
\code{average}  & no  & $6.047 \times 10^{-5}$ & $0.1234$ & $10.21$ & $0.2042$ \\
\code{average}  & yes & $5.702 \times 10^{-5}$ & $0.1199$ & $15.39$ & $0.1003$ \\
\code{ward}     & no  & $5.478 \times 10^{-5}$ & $0.1175$ & $13.27$ & $0.1377$ \\
\code{ward}     & yes & $5.397 \times 10^{-5}$ & $0.1166$ & $14.05$ & $0.1238$ \\
\bottomrule
\end{tabular}
\end{table}

Table~\ref{tab:linkage} shows the choice of linkage mattering roughly as much as
one would expect from the geometry. \code{single} linkage is the outlier: its
chaining produces the most concentrated allocation of the eight
($N_{\mathrm{eff}} = 7.82$ against $13.27$ for \code{ward}) and the highest
variance. This is the mechanism, not an accident --- an unbalanced tree splits a
budget unevenly by construction, because a deep leaf has been halved more often
than a shallow one.

Bisection improves every method on both measures here, and improves
\code{single} most (variance $-17\%$, $N_{\mathrm{eff}}$ from $7.82$ to
$13.63$), which is consistent with balance being what it supplies. We would not
read the uniform improvement as general: this is one panel, and the ordering of
two allocators on one covariance estimate is not evidence about their ordering
out of sample.

\subsection{The Schur parameter}

\begin{table}[t]
\centering
\caption{Schur complementary allocation on the same panel, \code{ward} linkage,
as $\gamma$ varies. The last row is the global minimum-variance portfolio,
which allows short positions, for reference.}
\label{tab:gamma}
\begin{tabular}{lrrrr}
\toprule
& Variance & Ann.\ vol. & $N_{\mathrm{eff}}$ & $\max_i w_i$ \\
\midrule
$\gamma = 0.00$ & $5.478 \times 10^{-5}$ & $0.1175$ & $13.27$ & $0.1377$ \\
$\gamma = 0.25$ & $5.473 \times 10^{-5}$ & $0.1174$ & $13.24$ & $0.1381$ \\
$\gamma = 0.50$ & $5.467 \times 10^{-5}$ & $0.1174$ & $13.20$ & $0.1385$ \\
$\gamma = 0.75$ & $5.459 \times 10^{-5}$ & $0.1173$ & $13.15$ & $0.1390$ \\
$\gamma = 1.00$ & $5.449 \times 10^{-5}$ & $0.1172$ & $13.08$ & $0.1397$ \\
\midrule
GMV (long/short) & $4.137 \times 10^{-5}$ & $0.1021$ & $4.37$ & $0.2438$ \\
\bottomrule
\end{tabular}
\end{table}

Table~\ref{tab:gamma} is the central empirical result of this paper, and it is
mostly a null one. Moving $\gamma$ from $0$ to $1$ --- from no cross-block
information to the full Schur complement --- reduces the portfolio variance by
$0.53\%$ and the annualised volatility by $3$ basis points. The allocation
barely moves: the largest weight changes by $0.002$, and $N_{\mathrm{eff}}$
falls from $13.27$ to $13.08$.

The reference row explains why that is worth stating. The global
minimum-variance portfolio on this covariance matrix has variance
$4.137 \times 10^{-5}$, a quarter below the $\gamma=1$ allocation, and it is a
visibly different object: $N_{\mathrm{eff}} = 4.37$, and it holds short
positions (its smallest weight is $-0.132$). Whatever $\gamma$ interpolates
towards in this implementation, on this panel it is not that portfolio.

\subsection{Three documented claims, examined}
\label{sec:claims}

The package documents three properties of the Schur allocator. We tested each.

\paragraph{$\gamma=0$ reproduces HRP exactly. Confirmed.}
The two allocators agree to the last bit on this panel: the maximum absolute
weight difference over the twenty assets is exactly $0$, not merely below a
tolerance. This is structural rather than fortunate --- at $\gamma=0$
Eq.~\eqref{eq:aug} leaves $A$ and $D$ untouched, so the two code paths compute
the same floating-point expression --- and it makes \code{schur\_hrp} a safe
drop-in for \code{hrp}.

\paragraph{$\gamma=1$ recovers the minimum-variance portfolio. Not as stated.}
Table~\ref{tab:gamma} already shows the two portfolios differing substantially
on the twenty-asset panel. The reason is visible in a case small enough to check
by hand. Take
\begin{equation}
  \Sigma = \begin{pmatrix} 4 & 1.5 \\ 1.5 & 1 \end{pmatrix},
\end{equation}
and the two-leaf tree. The global minimum-variance portfolio is
$w^{\star} = (-0.25,\, 1.25)$; the Schur allocator at $\gamma=1$ returns
$(0.2,\, 0.8)$.

The mechanism is Eq.~\eqref{eq:split}. Augmenting the blocks changes
\emph{which} variances are compared, but the comparison is still an
inverse-variance split, and an inverse-variance split of two sub-portfolios is
not the minimum-variance combination of them --- the minimum-variance weight for
two correlated blocks depends on their covariance, and no reweighting of the
diagonal blocks alone can reproduce it. A further consequence: because both
$\alpha$'s in Eq.~\eqref{eq:split} are non-negative, the output is long-only for
every $\gamma$, while the minimum-variance portfolio generally is not. The two
cannot coincide whenever $w^{\star}$ has a negative entry, whatever $\gamma$
does.

We are describing a gap between this implementation and the sentence in its own
documentation, not disputing \citet{cotton2024}; recovering minimum variance
requires the split itself to become a minimum-variance solve, which
Eq.~\eqref{eq:split} never becomes.

\paragraph{Variance falls as $\gamma$ rises. True here, not in general.}
It holds on the shipped panel, monotonically, as Table~\ref{tab:gamma} shows.
It is not a property of the algorithm. Over $200$ synthetic eight-asset panels
($60$ observations each, independent Gaussian log-returns; see
Appendix~\ref{app:repro}), the variance failed to be non-increasing in $\gamma$
in $76$ cases --- $38\%$. In the first such case the variance \emph{rose}
monotonically, from $1.0367 \times 10^{-5}$ at $\gamma=0$ to $1.0401
\times 10^{-5}$ at $\gamma=1$, an increase of $0.33\%$.

The effect is small in both directions, which is the practical point:
$\gamma$ is a weak knob on this construction. A user who reads
``raising $\gamma$ trades robustness for lower in-sample variance'' should not
expect the trade to be reliably available.

\section{Limitations}

The covariance estimator is the plain sample one. On a universe whose asset
count approaches its observation count this estimator is badly conditioned, and
while HRP never inverts it, the block variances of Eq.~\eqref{eq:blockvar} are
still computed from it and the Schur augmentation of Eq.~\eqref{eq:aug} does
solve against it. Shrinkage \citep{ledoit2004} is left to the caller, who may
pass any covariance frame.

The allocation is unconditionally long-only and fully invested. There is no
provision for weight caps, turnover limits, transaction costs, leverage, or a
risk-free asset, and no rebalancing logic: the package maps one covariance
estimate to one allocation. Nothing here evaluates HRP as a strategy. All
variances reported are in-sample, computed against the same covariance matrix
that produced the weights, and are properties of the construction rather than
evidence about performance --- the out-of-sample question that motivates HRP
\citep{lopezdeprado2016} is outside what this package measures.

\section{Conclusion}

Allocating on an explicit tree costs little and buys inspectability: the
intermediate budget of every cluster survives in the returned object, which
makes an HRP allocation auditable in a way a weight vector is not. The
construction is otherwise faithful to the published algorithm, and the two
variants the package adds behave as follows on the panel it ships. Bisection is
a real improvement in balance, and most valuable exactly where the linkage tree
is worst. The Schur parameter $\gamma$ is a much weaker instrument than its
documentation suggests: it moves the allocation by fractions of a percent, its
variance reduction is not monotone across panels, and its $\gamma=1$ endpoint is
not the minimum-variance portfolio, because the split rule it feeds stays
inverse-variance throughout. Users should treat $\gamma$ as a small perturbation
of HRP rather than as a dial between HRP and minimum variance, and the
documentation should say so.

\appendix

\section{Reproducibility}
\label{app:repro}

Every number in Section~\ref{sec:results} was produced with \code{pyhrp} 2.3.3
and the panel in the repository. Tables~\ref{tab:linkage} and~\ref{tab:gamma}:

\begin{verbatim}
import numpy as np, polars as pl
from pyhrp import hrp, schur_hrp, compute_returns, compute_cov

prices = pl.read_csv("tests/resources/stock_prices.csv",
                     try_parse_dates=True).drop("date")
cov = compute_cov(compute_returns(prices))

for method in ["single", "complete", "average", "ward"]:
    for bisection in [False, True]:
        root = hrp(prices=prices, method=method, bisection=bisection)
        w = np.array(list(root.portfolio.weights.values()))
        print(method, bisection, root.portfolio.variance(cov),
              1 / np.sum(w ** 2), w.max())

for gamma in [0.0, 0.25, 0.5, 0.75, 1.0]:
    root = schur_hrp(prices=prices, method="ward", gamma=gamma)
    print(gamma, root.portfolio.variance(cov))
\end{verbatim}

The minimum-variance reference row is
$w^{\star} \propto \Sigma^{-1} \mathbf{1}$, normalised to sum to one. The
monotonicity experiment of Section~\ref{sec:claims}:

\begin{verbatim}
rng = np.random.default_rng(0)
violations = 0
for _ in range(200):
    X = rng.standard_normal((60, 8))
    p = pl.DataFrame({f"A{i}": 100 * np.exp(np.cumsum(X[:, i] * 0.01))
                      for i in range(8)})
    c = compute_cov(compute_returns(p))
    v = [schur_hrp(prices=p, method="ward", gamma=g).portfolio.variance(c)
         for g in (0.0, 0.5, 1.0)]
    violations += not (v[0] >= v[1] >= v[2] - 1e-18)
print(violations, "/ 200")
\end{verbatim}

\bibliographystyle{plainnat}
\bibliography{references}

\end{document}
